\documentclass[a4paper,12pt]{article}




%\abovedisplayshortskip=-12pt plus 3pt
%\setlength{\abovedisplayskip}{-5pt}
%\setlength{\belowdisplayskip}{-5pt}
%\linespread{1.05}
%\voffset=-14cm
%\hoffset=-0.5cm

\usepackage{bm, amsmath, amsfonts, amssymb, url}
\usepackage{titlesec, titletoc}
\usepackage{graphicx,subfigure}
\usepackage{appendix, indentfirst, makeidx, array}
\usepackage{booktabs}
\usepackage{fancyhdr}
\pagestyle{fancy}
\usepackage{fullpage} 

\addtolength{\headsep}{+1cm}
\addtolength{\oddsidemargin}{+0.5cm}
\addtolength{\evensidemargin}{-0.5cm}
%\usepackage[cam,a4,center]{crop}
%\usepackage{geometry}


%\linespread{1.3}

%\fancypagestyle{plain}{
%\fancyhf{}
%\renewcommand{\headrulewidth}{0pt}
%%\fancyfoot[c]{\thepage}
%}


%\addtocontents{toc}{\protect\thispagestyle{empty}}

%%%-------page style ------------

%\setlength{\baselineskip}{20pt plus0.3pt minus0.2pt}
%\setlength{\marginparsep}{0pt}
%\setlength{\marginparwidth}{1.5cm}

%\setlength{\oddsidemargin}{3cm}
%\setlength{\evensidemargin}{3cm}
%%%-------------------------------


\renewcommand{\labelenumi}{(\alph{enumi})}
\renewcommand{\arraystretch}{1.2}
\titleformat{\chapter}[display]{\Large\bf}{\LARGE{Chapter~\thechapter}}{2mm}{}
\titleformat{\section}{\large\bf}{\thesection}{1.5mm}{}
\titleformat{\subsection}{\bfseries\itshape}{\thesubsection}{1.5mm}{}
%self-define symbol
\newcommand{\proposition}[1]{\bigskip\noindent\fbox{\begin{minipage}{0.95\textwidth}~\begin{minipage}{0.95\textwidth}#1\end{minipage}\\\smallskip\end{minipage}}\bigskip\vline\vline\vline}

\newcommand{\excercise}{\addtocounter{section}{+1}
\subsection*{\underline{Ex~\thesection}}}
\newcommand{\stag}[1]{\tag{\arabic{chapter}.\arabic{equation}#1}}

\newcommand{\qdo}[0]{\quad}
\newcommand{\qdt}[0]{\quad\quad}
\newcommand{\qdthr}[0]{\quad\quad\quad}
\newcommand{\qdf}[0]{\quad\quad\quad\quad}
\newcommand{\qdfiv}[0]{\quad\quad\quad\quad\quad}

\newcommand{\norm}[1]{\|#1\|}
%\newcommand{\sdelta}{{\scriptstyle\Delta}}
\newcommand{\sdelta}{\Delta}
\newcommand{\bra}[1]{\langle#1|}
\newcommand{\ket}[1]{|#1\rangle }
\newcommand{\bket}[1]{\Big |#1\Big\rangle }
\newcommand{\braket}[2]{\langle#1|#2\rangle }
\newcommand{\proj}[2]{|#1\rangle\langle#2|}
\newcommand{\expect}[2]{\langle #1\rangle_{#2}}
\newcommand{\Kappa}[0]{\bm{\mathcal{K}}}
\newcommand{\parity}[0]{\bm{\mathcal{P}}}
\newcommand{\parityl}[0]{\bm{p}}
\newcommand{\charge}[0]{\bm{\mathcal{C}}}


\makeindex
%-----------------------------------%

\begin{document}

\abovedisplayskip=-6pt plus 3pt minus 0pt
\belowdisplayskip=8pt plus 3pt minus 0pt

\abovedisplayshortskip=-6pt plus 3pt minus 0pt
\belowdisplayshortskip=8pt plus 3pt minus 0pt

\title{Problem Set 11 (Due: Dec. 14, 2011)}
\maketitle

Your homework should be submitted to TA (Chih-Han Lin) at the  beginning of the course. Generally, You have a week to finish the problem set since it was revealed. If you cannot submit it on time, you can still bring your homework to TA's office (L325, Institute of Atomic and Molecular Sciences, in NTU campus) and put it into TA's mailbox which is on the top of the shoe cabinet in L325. Always remember to remind TA to check your homework by e-mail (\url{clin@ltl.iams.sinica.edu.tw}) if you don't submit it in class.

%\abovedisplayskip=0pt %plus 3pt minus 3pt
%\belowdisplayskip=2pt %plus 3pt minus 3pt
%
%\abovedisplayshortskip=0pt %plus 3pt minus 3pt
%\belowdisplayshortskip=2pt %plus 3pt minus 3pt


\excercise 
Consider a 2$\times$2 matrix $\mathbf{A}$ for the linear transformation in 2-dimensional vector space, i.e.

$$
\begin{pmatrix}x'\\y'\end{pmatrix}=\mathbf{A}\begin{pmatrix}x\\y\end{pmatrix}
$$

\noindent such that the transformation leaves $x^2-y^2$ invariant, namely $x^2-y^2=x'^2-y'^2$.

\begin{enumerate}
\item Show that $\mathbf{A}$ satisfied

$$
\mathbf{A}^T\begin{pmatrix}1&0\\0&-1\end{pmatrix}\mathbf{A}=\begin{pmatrix}1&0\\0&-1\end{pmatrix}.
$$

\item Show that the transformation forms a group called O(1,1)-group.

\item Construct explicitly the matrix $\mathbf{A}$ with det $\mathbf{A}=1$.

\item Construct explicitly the generator of $\mathbf{A}$ and reconstruct $\mathbf{A}$ through its generator.
\end{enumerate}


\excercise

Consider the following transformation,

$$
x'=a^1x+a^2.
$$

\begin{enumerate}
\item What are the conditions if the transformation forms a 2 parameters group?
\item Denote the group element by $ a=(a^1,a^2) $. Find the inverse element of $ a $, namely, $ a^{-1}=((a^{-1})^1,(a^{-1})^2) $.
\item What are the composition rules of the group parameters, i.e. try to find $ c=ba $ for $ c^l=c^l(a^i,b^j)=f^l(a^i,b^j) $?
\item Can you construct the $ 2\times 2 $ matrix representation of the group to justify
your answer in (a), (b) and (c)?
\end{enumerate}

\excercise

Consider the rotational transformation about the azimuthal axis with an angle $ \theta $, namely,

\begin{align*}
x'^1&=x^1\cos\theta-x^2\sin\theta,\\
x'^2&=x^1\sin\theta+x^2\cos\theta,\\
x'^3&=x^3,
\end{align*}

\noindent then the function $ F(x)=F(x^1,x^2,x^3) $ will be changed into 

$$
F(x')=F(x^1\cos\theta-x^2\sin\theta,x^1\sin\theta+x^2\cos\theta,x^3).
$$

Prove explicitly that the transformed function can be obtained by

\begin{align*}
\exp(\theta X_3)F(x)&=\exp\left[\theta\left(-x^2\frac{\partial}{\partial x^1}+x^1\frac{\partial}{\partial x^2}\right)\right]F(x^1,x^2,x^3)\\
&=F(x'^1,x'^2,x'^3).
\end{align*}

(\textbf{Hint:} use $ \exp(\theta X_3)=\exp\left(\theta\partial/\partial\varphi\right) $ and write 

$$
F(x) = F(r\cos\varphi,r\sin\varphi, x^3).)
$$




\end{document}
